\item Los astronautas del Apollo 11 colocaron un panel de retrorreflectores de esquinas cúbicas eficientes en la superficie de la Luna. La velocidad de la luz se deduce al medir el intervalo de tiempo necesario para que un láser se dirija desde la Tierra, se refleje en el panel y regrese a la Tierra. Suponga que la medición de este intervalo es $2.51\text{ s}$ en la estación cuando la Luna está en su cenit. Considere la distancia de centro a centro de la Tierra a la Luna en $3.84 \times 10^8\text{ m}$.
\begin{enumerate}
    \item ¿Cuál es la velocidad medida de la luz?
    \item Explique si es necesario tomar en cuenta los tamaños de la Tierra y de la Luna en sus cálculos.
\end{enumerate}